% !TEX root =  Lecture8.tex


%%%%%%%%%%%%%%%%%%%%%%%%%%%%%%%%%%%%
% Recap/Agenda
%%%%%%%%%%%%%%%%%%%%%%%%%%%%%%%%%%%%
\begin{frame}
\frametitle{Module 2: Statistical Inference}
    \begin{itemize}
    \item \hl{Previously:} Module 1 --- building, interpreting, and choosing regression models.
    \item \hl{This time:} \emph{how sure are we?} Testing with randomization.
      \begin{itemize}
      \item statistics, sampling distributions, and the null distribution
      \item hypothesis testing review, p-values, and decision errors
      \item building the null distribution by \hl{randomization}
      \end{itemize}
    \item \hl{Reading:} IMS Chapters 11 and 14.
    \item \hl{Homework for today:}
    \item \hl{Homework for next class:}
    \end{itemize}
\end{frame}


%%%%%%%%%%%%%%%%%%%%%%%%%%%%%%%%%%%%
% Learning objectives:
%%%%%%%%%%%%%%%%%%%%%%%%%%%%%%%%%%%%
\begin{frame}
    \frametitle{Lecture Objectives}
    \goals{%
    \begin{itemize}
    \item explain what a sampling distribution is, and what makes the \emph{null} distribution the one we can build;
    \item state hypotheses, choose a test statistic, and interpret a p-value correctly;
    \item distinguish Type I from Type II errors and say which one matters in context;
    \item given a null hypothesis, determine \emph{what to randomize} --- and carry out the test.
    \end{itemize}}
    \objectives{%
    \begin{itemize}
    \item \textbf{(M2, LO1)} Approaches to Inference
    \item \textbf{(M2, LO2)} Simulation-based Inference
    \end{itemize}}
\end{frame}
